\section{Mô phỏng thuật toán}

\subsection{Mục tiêu mô phỏng}

Mục tiêu của phần mô phỏng là kiểm tra hành vi thực nghiệm của các thuật toán tối ưu dưới mô-đun trên hai bài toán mẫu: bài toán phủ sóng tối đa và bài toán lựa chọn tập đại diện. Các thí nghiệm tập trung vào giá trị hàm mục tiêu, tỉ lệ so với nghiệm tối ưu trong trường hợp dữ liệu nhỏ, số lần đánh giá hàm mục tiêu hoặc lợi ích biên, và thời gian chạy.

\subsection{Thiết lập thực nghiệm}

Đối với bài toán phủ sóng, thí nghiệm sử dụng năm tập con \(A_1,\ldots,A_5\) và ràng buộc chọn \(k=2\). Các thuật toán được so sánh gồm brute force, greedy và random baseline.

Đối với bài toán lựa chọn tập đại diện, dữ liệu nhỏ gồm hai cụm Gaussian, mỗi cụm có 5 điểm, với ràng buộc chọn \(k=2\). Độ tương tự giữa hai điểm được tính bởi hàm Gaussian với \(\sigma=1.0\). Các thuật toán được so sánh trên dữ liệu nhỏ gồm brute force, greedy và random baseline.

Đối với thí nghiệm thời gian chạy, dữ liệu gồm hai cụm Gaussian với số điểm lần lượt là \(n=200\), \(n=500\) và \(n=1000\). Ràng buộc chọn là \(k=10\), \(\sigma=1.0\). Các thuật toán được so sánh gồm greedy, lazy greedy và stochastic greedy với \(\epsilon=0.1\).

\subsection{Kết quả trên bài toán phủ sóng}

Bảng~\ref{tab:max-coverage-small} trình bày kết quả trên bài toán phủ sóng nhỏ.

\begin{table}[ht]
    \centering
    \caption{Kết quả mô phỏng trên bài toán phủ sóng nhỏ}
    \label{tab:max-coverage-small}
    \begin{tabular}{llllll}
\hline
Algorithm & Selected & Objective Value & Ratio to Optimal & Evaluations & Runtime Seconds \\
\hline
Brute Force & {A1, A4} & 6.000000 & 1.000000 & 10 & 0.000018 \\
Greedy & {A1, A4} & 6.000000 & 1.000000 & 10 & 0.000015 \\
Random Baseline & {A2, A4} & 6.000000 & 1.000000 & 1000 & 0.002853 \\
\hline
\end{tabular}

\end{table}

Theo kết quả trong bảng, brute force chọn \(\{A1,A4\}\) và đạt giá trị hàm mục tiêu \(6.0\). Greedy cũng chọn \(\{A1,A4\}\), đạt giá trị \(6.0\), tương ứng tỉ lệ \(1.0\) so với nghiệm tối ưu. Random baseline chọn \(\{A2,A4\}\), cũng đạt giá trị \(6.0\) trong lần chạy này, nhưng dùng \(1000\) lần đánh giá.

\subsection{Kết quả trên bài toán lựa chọn tập đại diện}

Bảng~\ref{tab:facility-location-small} trình bày kết quả trên bài toán lựa chọn tập đại diện với dữ liệu nhỏ.

\begin{table}[ht]
    \centering
    \caption{Kết quả mô phỏng trên bài toán lựa chọn tập đại diện nhỏ}
    \label{tab:facility-location-small}
    \begin{tabular}{llllll}
\hline
Algorithm & Selected & Objective Value & Ratio to Optimal & Evaluations & Runtime Seconds \\
\hline
Brute Force & {3, 9} & 9.089966 & 1.000000 & 45 & 0.001048 \\
Greedy & {3, 9} & 9.089966 & 1.000000 & 20 & 0.000690 \\
Random Baseline & {3, 9} & 9.089966 & 1.000000 & 1000 & 0.020225 \\
\hline
\end{tabular}

\end{table}

Trên dữ liệu nhỏ, brute force và greedy đều chọn tập chỉ số \(\{3,9\}\), đạt giá trị hàm mục tiêu \(9.089966054828349\) và tỉ lệ \(1.0\) so với nghiệm tối ưu. Random baseline trong lần chạy này cũng chọn \(\{3,9\}\), đạt cùng giá trị mục tiêu, nhưng cần \(1000\) lần đánh giá.

\begin{figure}[ht]
    \centering
    \includegraphics[width=0.75\textwidth]{Hinhve/facility_location_small.png}
    \caption{Dữ liệu hai cụm Gaussian và các điểm đại diện được greedy chọn}
    \label{fig:facility-location-small}
\end{figure}

\subsection{So sánh thời gian chạy}

Bảng~\ref{tab:runtime-comparison} trình bày kết quả so sánh thời gian chạy và số lần đánh giá trên bài toán lựa chọn tập đại diện với các quy mô dữ liệu khác nhau.

\begin{table}[ht]
    \centering
    \caption{So sánh thời gian chạy và số lần đánh giá}
    \label{tab:runtime-comparison}
    \begin{tabular}{lllll}
\hline
n & Algorithm & Objective Value & Evaluations & Runtime Seconds \\
\hline
200 & Greedy & 191.750321 & 1956 & 0.035055 \\
200 & Lazy Greedy & 191.750321 & 638 & 0.011534 \\
200 & Stochastic Greedy & 191.484586 & 471 & 0.008990 \\
500 & Greedy & 476.635032 & 4956 & 0.107231 \\
500 & Lazy Greedy & 476.635032 & 1664 & 0.035912 \\
500 & Stochastic Greedy & 476.452069 & 1161 & 0.026410 \\
1000 & Greedy & 952.167796 & 9956 & 0.339457 \\
1000 & Lazy Greedy & 952.167796 & 3130 & 0.122480 \\
1000 & Stochastic Greedy & 950.790369 & 2311 & 0.082870 \\
\hline
\end{tabular}

\end{table}

Với \(n=200\), greedy đạt giá trị \(191.75032124917044\) với \(1956\) lần đánh giá và thời gian \(0.036216100001183804\) giây. Lazy greedy đạt cùng giá trị mục tiêu với \(638\) lần đánh giá và \(0.01186020000022836\) giây. Stochastic greedy đạt giá trị \(191.48458574287434\), dùng \(471\) lần đánh giá và \(0.009123300002102042\) giây.

Với \(n=500\), greedy và lazy greedy cùng đạt giá trị \(476.6350317731388\). Số lần đánh giá của greedy là \(4956\), trong khi lazy greedy là \(1664\). Stochastic greedy đạt giá trị \(476.4520694037594\) với \(1161\) lần đánh giá.

Với \(n=1000\), greedy và lazy greedy cùng đạt giá trị \(952.1677964048943\). Greedy dùng \(9956\) lần đánh giá và \(0.35698779999802355\) giây, còn lazy greedy dùng \(3130\) lần đánh giá và \(0.10470129999885103\) giây. Stochastic greedy đạt giá trị \(950.7903692993278\) với \(2311\) lần đánh giá và \(0.08070120000047609\) giây.

\begin{figure}[ht]
    \centering
    \includegraphics[width=0.75\textwidth]{Hinhve/runtime_comparison.png}
    \caption{So sánh thời gian chạy theo quy mô dữ liệu}
    \label{fig:runtime-comparison}
\end{figure}

\begin{figure}[ht]
    \centering
    \includegraphics[width=0.75\textwidth]{Hinhve/evaluation_comparison.png}
    \caption{So sánh số lần đánh giá theo quy mô dữ liệu}
    \label{fig:evaluation-comparison}
\end{figure}

\subsection{Nhận xét}

Trên các dữ liệu nhỏ, greedy đạt cùng giá trị mục tiêu với brute force trong cả hai bài toán được xét. Điều này phù hợp với mục tiêu của thí nghiệm nhỏ: dùng brute force làm mốc so sánh và kiểm tra kết quả của thuật toán xấp xỉ.

Trong thí nghiệm thời gian chạy, lazy greedy đạt cùng giá trị mục tiêu với greedy ở cả ba quy mô dữ liệu, đồng thời giảm số lần đánh giá từ \(1956\), \(4956\), \(9956\) xuống còn \(638\), \(1664\), \(3130\). Stochastic greedy có số lần đánh giá và thời gian chạy nhỏ hơn greedy, nhưng giá trị hàm mục tiêu thấp hơn nhẹ trong các lần chạy được ghi nhận.
